% Figure: convergence of the Fourier sine series for a smooth (sine)
% datum versus a discontinuous (square) datum. Plots the L2 partial-sum
% error against the number of retained terms on log-log axes, showing
% the smoothness-controlled convergence rate.
\begin{figure}[htbp]
\centering
\begin{tikzpicture}
\begin{loglogaxis}[
  width=12cm, height=6.5cm,
  xlabel={number of retained terms $N$},
  ylabel={$L^2$ truncation error $\|f - S_N\|_2$},
  xmin=1, xmax=200,
  ymin=1e-4, ymax=1,
  grid=both,
  major grid style={engblack!20},
  minor grid style={engblack!10},
  legend pos=south west,
  legend cell align=left,
  font=\small,
]
% Square wave on [0,L]: b_n ~ 1/n for odd n; tail energy beyond N
% ~ sum_{n>N, odd} 1/n^2 ~ 1/(2N). So L2 error ~ C / sqrt(N).
% Use error = 0.45 / sqrt(N).
\addplot[engred, very thick, mark=square*, mark size=1.3pt,
         domain=1:200, samples=18] {0.45/sqrt(x)};
\addlegendentry{square datum: error $\sim N^{-1/2}$}

% Triangle (continuous, kinked) datum: b_n ~ 1/n^2, tail energy
% beyond N ~ sum_{n>N} 1/n^4 ~ 1/(3 N^3). L2 error ~ C N^{-3/2}.
\addplot[engblue, very thick, mark=*, mark size=1.3pt,
         domain=1:200, samples=18] {0.30/(x*sqrt(x))};
\addlegendentry{triangle datum: error $\sim N^{-3/2}$}

% reference slopes
\addplot[engblack, dashed, domain=1:200, samples=2] {0.45/sqrt(x)};
\addplot[engblack, dotted, domain=1:200, samples=2] {0.30/(x*sqrt(x))};
\end{loglogaxis}
\end{tikzpicture}
\caption[Fourier convergence rate set by smoothness]{$L^2$
truncation error of the Fourier sine series against the number of
retained terms $N$, on log-log axes. The square-wave datum (red) has
a jump discontinuity; its coefficients fall as $1/n$ and the
truncation error decays only as $N^{-1/2}$ (slope $-1/2$). The
triangle datum (blue) is continuous with a kinked derivative; its
coefficients fall as $1/n^2$ and the truncation error decays as
$N^{-3/2}$ (slope $-3/2$). Each extra order of smoothness of the
datum buys one extra power of $n$ in the coefficient decay and one
extra half-power in the $L^2$ convergence rate. The pointwise
behaviour at the jump is the Gibbs overshoot
(\cref{fig:vol02:ch12:gibbs}), which is invisible in the $L^2$ norm
because its width shrinks even as its height does not. The working
moral: read the convergence rate off the smoothness of the data
before counting terms.}
\label{fig:vol02:ch12:fourier-convergence}
\end{figure}
